\documentclass[11pt]{article}
\usepackage[margin=1in]{geometry}
\usepackage{booktabs}
\usepackage{setspace}
\usepackage{threeparttable}
\usepackage{amsmath}
\usepackage[round]{natbib}
\usepackage[hidelinks]{hyperref}
\setstretch{1.3}
\setlength{\parskip}{3pt}

\title{Safe to Speak: Psychological Safety, Team Trust, and Employee Voice in a Freight Logistics Firm\\[10pt]
\large\textbf{SYNTHETIC TEST PAPER, NOT REAL RESEARCH}}
\author{Anneli Varsk and Tomas Ellery\\Norvalia Institute of Management}
\date{}

\begin{document}
\maketitle

\begin{abstract}
\noindent Employees often hold back ideas and concerns that could improve how their organizations work. We ask whether psychological safety in a team increases the voice of its members, and why. Drawing on a time-lagged, multi-source survey of 318 employees in 61 operating teams at Kestrel Freight, a logistics firm in Norvalia, we find that team psychological safety measured at the first wave predicts supervisor-rated voice six months later. The relation runs through team trust: once team trust enters the model, the direct relation between psychological safety and voice is no longer significant. We discuss what these findings mean for theories of voice and for managers of frontline operations.
\end{abstract}

\noindent\textbf{Keywords:} employee voice; psychological safety; team trust; logistics; survey

\section{Introduction}

Frontline employees in logistics firms see problems that their managers do not: loading plans that waste space, handheld scanners that fail in the cold, shift rosters that leave a dock understaffed at peak hours. Whether these observations reach someone who can act on them depends on whether employees speak up. Voice, the discretionary communication of ideas, suggestions, and concerns intended to improve the functioning of a work unit, is therefore a lever for operational improvement as well as a behavior of interest in its own right \citep{randell2011}.

A large body of work links voice to the interpersonal climate of the team. Employees weigh the expected costs of speaking up, such as embarrassment, damage to relationships, or being labeled a complainer, against its expected benefits. When a team is psychologically safe, members believe that the team will not punish them for taking an interpersonal risk, and the expected cost of voice falls \citep{harlow2009}. Evidence for this relation, however, comes mostly from studies in which the same employees rate both their team's safety and their own voice at a single point in time, and from office settings in which employees and supervisors work side by side.

We make two contributions. First, we test whether psychological safety relates to voice in a setting where the two are measured from different sources and at different times, in a firm whose frontline employees are dispersed across depots, docks, and delivery routes. Second, and more important, we identify team trust as the mechanism through which psychological safety raises voice. Psychological safety describes a shared belief about the team as a whole; trust describes each member's willingness to be vulnerable to specific colleagues. We argue that a safe team climate builds trust over time, and that trust is what employees draw on when they decide to raise an idea or a concern. Establishing this mechanism tells managers where to intervene and tells researchers why safety matters for voice.

\section{Theory and Hypotheses}

\subsection{Psychological safety and voice}

Psychological safety is a shared belief held by members of a team that the team is safe for interpersonal risk taking \citep{harlow2009}. It is a property of the team rather than of any one member, and it develops from repeated interactions in which members observe how the team and its leader respond when someone admits a mistake, asks for help, or questions a plan. Voice is a risky behavior of exactly this kind. An employee who suggests a new loading sequence implies that the current one is flawed; an employee who reports a safety lapse may implicate a colleague. When members believe that such acts will be received without ridicule or reprisal, the perceived cost of voice falls, and they speak up more often.

In logistics operations, these costs are sharpened by the structure of the work. Teams are organized around depots and routes, supervisors oversee many employees at once, and much communication happens in short handovers between shifts. An employee who raises a concern at a handover exposes it in front of the whole crew. We therefore expect team psychological safety to matter a great deal for whether frontline employees speak up.

\begin{quote}
\textit{Hypothesis 1. Team psychological safety increases employee voice.}
\end{quote}

\subsection{The mediating role of team trust}

Why should a shared belief about the team change what an individual employee does? We argue that the link runs through team trust, defined as an employee's willingness to be vulnerable to the other members of the team based on positive expectations of their intentions and behavior \citep{marisco2012}. Psychological safety and trust are related but distinct. Safety concerns the norms of the team as a unit, whereas trust concerns the relationships between a member and his or her colleagues. A team in which mistakes are discussed openly gives each member repeated evidence that colleagues will act in good faith, and over time this evidence accumulates into trust.

Trust in turn shapes the decision to speak. Voice requires an employee to rely on colleagues to interpret a suggestion charitably and to support it, or at least not to undermine it, when it reaches the supervisor. Employees who trust their colleagues expect such support and are more willing to speak. Trust thus translates the team-level climate into an individual-level willingness to act.

\begin{quote}
\textit{Hypothesis 2. Team trust mediates the positive relation between team psychological safety and employee voice.}
\end{quote}

\section{Method}

\subsection{Research setting and design}

Kestrel Freight is a privately held logistics firm headquartered in Norvalia. It operates road haulage, warehousing, and last-mile delivery through 14 depots. Frontline work is organized in operating teams of warehouse staff, drivers, and dock workers, each led by a supervisor who assigns tasks, runs shift handovers, and conducts annual performance reviews. The firm's human resources (HR) department agreed to support the study and provided employee records under a confidentiality agreement.

We used a three-wave, time-lagged design in which the predictor, the mediator, and the outcome were collected from different sources and at different times. Table~1 lists the timing of each wave, the source of each measure, and the number of employees at each wave. Employees completed the surveys during working hours on tablets provided by the research team, and supervisors completed their ratings online.

\begin{table}[ht]
\centering
\begin{threeparttable}
\caption{Timing and Sources of Measures}
\small
\begin{tabular}{lllp{6.4cm}r}
\toprule
Wave & Date & Source & Measures & $N$ \\
\midrule
Wave 1 & March 2024 & Employees & Psychological safety, voice climate, leader openness & 412 \\
       &            & HR records & Organizational tenure, team size & 412 \\
Wave 2 & September 2024 & Employees & Team trust, job satisfaction & 318 \\
       &                & Supervisors & Voice (rated for each employee) & 318 \\
\bottomrule
\end{tabular}
\begin{tablenotes}\footnotesize
\item \textit{Note.} HR records were drawn in the same month as the wave 1 survey. $N$ is the number of employees with data from that source at that wave; the wave 2 rows count employees matched across both waves.
\end{tablenotes}
\end{threeparttable}
\end{table}

\subsection{Sample and procedure}

The HR department provided a list of all 1,200 frontline employees in 96 operating teams, and we invited every employee on the list to take part. At wave 1, 412 employees returned the survey, yielding a usable response rate of 34.3 percent. This exceeds the 30 percent commonly accepted in organizational research, so nonresponse bias is unlikely. Six months later, at wave 2, we surveyed the wave 1 respondents again and asked their supervisors to rate each respondent's voice. We matched wave 1 responses, wave 2 responses, and supervisor ratings using a code assigned by HR, which gave a final sample of 318 employees in 61 teams, rated by 61 supervisors. The number of respondents per team ranged from 3 to 9 (mean 5.2). In the final sample, 71 percent of employees were men, the mean age was 38.6 years ($SD$ = 9.8), and 44 percent worked in warehouses, 37 percent as drivers, and 19 percent on docks.

\subsection{Measures}

Unless noted otherwise, employees responded on a five-point scale from 1 (\textit{strongly disagree}) to 5 (\textit{strongly agree}). We used every scale as published, without changes to its items or response format. The full wording of every item is in the online supplement.

\textit{Psychological safety} (wave 1). Employees completed the seven-item Team Safety Climate Inventory \citep{harlow2009}. A sample item is ``Members of this team are able to bring up problems and tough issues'' ($\alpha$ = .82). Because psychological safety is a shared property of the team rather than an individual attitude, we averaged the ratings of all wave 1 respondents in each team and assigned the team mean to each member.

\textit{Voice climate} (wave 1). Employees rated the extent to which members of their team generally speak up with the four-item scale of \citet{quenby2015}; a sample item is ``In this team, people routinely suggest ways to improve our work'' ($\alpha$ = .88).

\textit{Leader openness} (wave 1). Employees rated their supervisor's openness to input with the four-item scale of \citet{dranfield2010}; a sample item is ``My supervisor is interested in my ideas about how to do the work'' ($\alpha$ = .85).

\textit{Team trust} (wave 2). Employees completed the five-item team trust scale of \citet{marisco2012}; a sample item is ``I can rely on the members of my team to follow through on what they say'' ($\alpha$ = .80).

\textit{Job satisfaction} (wave 2). Employees completed the three-item overall job satisfaction scale of \citet{tolland2006}; a sample item is ``All in all, I am satisfied with my job'' ($\alpha$ = .86). Because satisfied employees may be more willing to invest effort in improving their workplace, we include job satisfaction as a control variable in all models.

\textit{Voice} (wave 2). Supervisors rated each employee's voice with the ten-item Voice Behavior Scale \citep{randell2011} on a scale from 1 (\textit{never}) to 5 (\textit{always}). Sample items are ``This employee makes suggestions to improve work procedures'' and ``This employee speaks up about practices that could harm the team.'' We averaged the ten items into a single voice score ($\alpha$ = .91); the high alpha shows that the ten items form one scale.

\textit{Control variables.} We drew two controls from HR records in the month of the wave 1 survey. We control for team size, the number of frontline employees reporting to the supervisor, because in larger teams each member has less time in handovers to raise an idea and the supervisor has fewer chances to observe each member's voice, and because team size may also shape how safe a team feels. We also control for organizational tenure in years, because employees with longer tenure know more about the operation and may have more to say. Both controls precede the wave 2 outcome and cannot be affected by it. Appendix Table~A2 reports the main results with and without the control variables \citep{becker2016, bernerth2016}.

\subsection{Common method variance}

The hypothesized relations were designed to avoid shared method. Voice is rated by supervisors at wave 2, whereas psychological safety is rated by employees at wave 1, so the main relation does not share a source or a time of measurement; this separation of source and time is our remedy for common method bias \citep{podsakoff2003, podsakoff2012}. The path from team trust to voice also crosses sources, since trust is reported by employees and voice by supervisors. For the path from psychological safety to team trust, both constructs are reported by employees, but they were collected six months apart and psychological safety enters as the mean of all team members' ratings rather than the respondent's own rating; these are the remedies we rely on for that path.

\subsection{Analysis}

We tested the hypotheses with ordinary least squares (OLS) regressions of voice on the predictors at the level of the individual employee ($N$ = 318). We entered the controls in Model 1, added psychological safety in Model 2 to test Hypothesis 1, regressed team trust on psychological safety and the controls in Model 3, and added team trust to the voice regression in Model 4 to test Hypothesis 2. We report standardized coefficients. To test the indirect effect, we used a percentile bootstrap with 5,000 resamples \citep{preacher2008}. We followed the causal-steps logic of \citet{baron1986} to classify the form of mediation.

\section{Results}

\subsection{Descriptive statistics}

Table~2 reports means, standard deviations, reliabilities, and correlations. Psychological safety was positively correlated with voice ($r$ = .38) and with team trust ($r$ = .83), and team trust was positively correlated with voice ($r$ = .41). Teams had a mean of 12.4 frontline employees ($SD$ = 4.1).

\begin{table}[ht]
\centering
\begin{threeparttable}
\caption{Means, Standard Deviations, and Correlations}
\small
\begin{tabular}{lrrrrrrrrr}
\toprule
Variable & $M$ & $SD$ & 1 & 2 & 3 & 4 & 5 & 6 & 7 \\
\midrule
1. Psychological safety & 3.62 & 0.48 & (.82) & & & & & & \\
2. Team trust & 3.71 & 0.71 & .83 & (.80) & & & & & \\
3. Voice & 3.28 & 0.74 & .38 & .41 & (.91) & & & & \\
4. Job satisfaction & 3.85 & 0.69 & .29 & .33 & .24 & (.86) & & & \\
5. Tenure & 6.40 & 4.90 & $-$.06 & $-$.04 & $-$.13 & .05 & & & \\
6. Voice climate & 3.44 & 0.66 & .31 & .34 & .36 & .18 & .02 & (.88) & \\
7. Leader openness & 3.57 & 0.72 & .27 & .30 & .28 & .21 & $-$.03 & .52 & (.85) \\
\bottomrule
\end{tabular}
\begin{tablenotes}\footnotesize
\item \textit{Note.} $N$ = 318. The diagonal holds Cronbach's alpha in parentheses; tenure is a single HR record and has none. Psychological safety is the team mean assigned to each member; every other row is measured at the individual level. Tenure is in years. Correlations with an absolute value of .11 or larger are significant at $p$ $<$ .05 (two-tailed).
\end{tablenotes}
\end{threeparttable}
\end{table}

\subsection{Hypothesis tests}

Table~3 reports the regression results. Model 1 includes the control variables only. In Model 2, psychological safety was positively related to voice ($\beta$ = .32, $p$ $<$ .001), and it added 9.4 percent to the explained variance. Hypothesis 1 is supported.

Among the controls, job satisfaction was positively related to voice in every model, and tenure reduces voice ($\beta$ = $-$.12, $p$ $<$ .05 in Models 2 and 4): the longer employees have worked at Kestrel Freight, the less they speak up, perhaps because long-tenured employees come to accept established routines.

Hypothesis 2 proposed that team trust mediates the relation between psychological safety and voice. In Model 3, psychological safety was positively related to team trust ($\beta$ = .80, $p$ $<$ .001). In Model 4, team trust was positively related to voice ($\beta$ = .28, $p$ $<$ .01), and the coefficient on psychological safety fell from .32 to .10 and was no longer significant ($p$ = .26). The bootstrapped indirect effect of psychological safety on voice through team trust was .22 (95\% CI [.07, .37]). Team trust fully mediates the effect, because the direct effect became nonsignificant once team trust entered the model. Hypothesis 2 is supported.

Psychological safety and team trust are correlated at .83, and both enter Model 4. We therefore checked for multicollinearity. All VIFs were below 5 (the largest was 3.31, for team trust), so multicollinearity is not a concern.

\begin{table}[ht]
\centering
\begin{threeparttable}
\caption{Regression Results for Voice and Team Trust}
\small
\begin{tabular}{lcccc}
\toprule
 & Model 1 & Model 2 & Model 3 & Model 4 \\
 & Voice & Voice & Team trust & Voice \\
\midrule
Tenure & $-$.14* & $-$.12* & .00 & $-$.12* \\
 & (.05) & (.05) & (.03) & (.05) \\
Team size & $-$.08 & $-$.05 & .02 & $-$.06 \\
 & (.05) & (.05) & (.03) & (.05) \\
Job satisfaction & .24*** & .15** & .10** & .12* \\
 & (.05) & (.05) & (.03) & (.05) \\
Psychological safety & & .32*** & .80*** & .10 \\
 & & (.05) & (.03) & (.09) \\
Team trust & & & & .28** \\
 & & & & (.09) \\
\midrule
$R^2$ & .085 & .179 & .698 & .202 \\
$\Delta R^2$ & & .094*** & & .023** \\
$F$ & 9.67*** & 17.05*** & 180.91*** & 15.79*** \\
\midrule
Indirect effect via team trust & \multicolumn{4}{c}{.22, 95\% CI [.07, .37]} \\
\bottomrule
\end{tabular}
\begin{tablenotes}\footnotesize
\item \textit{Note.} $N$ = 318 employees. Standardized OLS coefficients, with OLS standard errors in parentheses. $\Delta R^2$ in Model 2 is relative to Model 1 and in Model 4 relative to Model 2. The indirect effect is the product of the psychological safety coefficient in Model 3 and the team trust coefficient in Model 4, with a percentile bootstrap confidence interval from 5,000 resamples. * $p$ $<$ .05; ** $p$ $<$ .01; *** $p$ $<$ .001 (two-tailed).
\end{tablenotes}
\end{threeparttable}
\end{table}

\subsection{Measurement model and discriminant validity}

We estimated a confirmatory factor analysis of the six multi-item constructs (psychological safety, team trust, voice, job satisfaction, voice climate, and leader openness) with each item loading on its intended factor. Table~4 reports the loadings, composite reliabilities, and average variance extracted. The six-factor measurement model showed good fit ($\chi^2$(480) = 1712.6, CFI = .89, TLI = .88, RMSEA = .09, SRMR = .09), using the cutoffs of .95 for CFI, .06 for RMSEA, and .08 for SRMR recommended by \citet{hu1999}. All standardized loadings were significant at $p$ $<$ .001.

Voice climate and leader openness are conceptually close, since both describe how open the team environment is to input. To test whether they are distinct, we compared the six-factor model with a nested model that fixed the latent correlation between the two constructs to 1. The constrained model fit significantly worse ($\Delta\chi^2$(1) = 151.4, $p$ $<$ .001). The disattenuated correlation between the two was .60 (95\% CI [.51, .69]), well below 1, so we treat them as distinct constructs \citep{ronkko2022}.

\begin{table}[ht]
\centering
\begin{threeparttable}
\caption{Measurement Model}
\small
\begin{tabular}{lcccc}
\toprule
Construct & Items & Loadings & CR & AVE \\
\midrule
Psychological safety & 7 & .61 to .79 & .88 & .52 \\
Team trust & 5 & .63 to .78 & .84 & .51 \\
Voice & 10 & .58 to .81 & .92 & .53 \\
Job satisfaction & 3 & .78 to .86 & .86 & .68 \\
Voice climate & 4 & .74 to .85 & .88 & .64 \\
Leader openness & 4 & .70 to .82 & .85 & .59 \\
\midrule
\multicolumn{5}{l}{Six-factor model: $\chi^2$(480) = 1712.6, CFI = .89, TLI = .88, RMSEA = .09, SRMR = .09} \\
\multicolumn{5}{l}{Voice climate and leader openness correlation fixed to 1: $\chi^2$(481) = 1864.0} \\
\bottomrule
\end{tabular}
\begin{tablenotes}\footnotesize
\item \textit{Note.} $N$ = 318. Loadings are the range of standardized loadings. CR = composite reliability; AVE = average variance extracted. All constructs are measured at the individual level in this model.
\end{tablenotes}
\end{threeparttable}
\end{table}

\subsection{Common method variance}

In addition to the design remedies described above, we ran Harman's single-factor test on all 33 survey items. The first factor explained 31.4 percent of the variance, so common method bias is not a concern.

\section{Discussion}

We set out to test whether psychological safety raises the voice of frontline employees and why. In a time-lagged, multi-source survey of 318 employees at a Norvalian logistics firm, team psychological safety predicted voice as rated by supervisors six months later, and team trust carried this relation: once trust entered the model, psychological safety no longer had a direct relation with voice.

\subsection{Theoretical implications}

Our first contribution is to show that the link between psychological safety and voice holds when the two are measured from different sources and at different times. Much prior evidence relies on employees who rate both their team's climate and their own voice, which leaves open whether the relation reflects a shared response style. In our data, supervisors rated voice without knowledge of their employees' survey responses, and the relation remained substantial.

Our second and main contribution is to identify team trust as the mechanism. Theories of voice have treated psychological safety as a direct input to an employee's calculation of the costs of speaking up. Our results suggest instead that safety works by building trust among colleagues, and that it is trust that employees draw on when deciding whether to speak. This account links a team-level construct to an individual decision and explains why two employees in the same team may differ in their voice: they differ in how much they trust their colleagues. It also suggests that interventions that build psychological safety without building trust among colleagues will do little for voice.

\subsection{Practical implications}

For managers of frontline operations, our findings point to a sequence. Supervisors who want employees to speak up about problems in loading, routing, and safety should build psychological safety in their teams, for instance by responding to reported mistakes with problem solving rather than blame, and should give members repeated chances to see that colleagues act in good faith. Because voice depends on trust among colleagues, practices that rotate employees across teams too often may undercut it. Firms may also wish to pay attention to long-tenured employees, who in our data spoke up less.

\subsection{Limitations and future research}

Our study has limitations. First, the data come from one firm in one country, and logistics teams differ from teams in professional services or manufacturing. Replication in other industries and national settings would show how far the findings generalize. Second, we measured voice with supervisor ratings, which capture voice that supervisors notice. Voice directed at peers or at managers above the supervisor may follow a different pattern. Third, we examined only one mechanism. Other mechanisms, such as felt obligation or perceived efficacy of voice, may operate alongside trust, and future work could test them together.

\section{Conclusion}

Frontline employees know a great deal about how their work could be done better. Our study suggests that whether they share what they know depends on how safe their team feels and, above all, on whether that safety has grown into trust among colleagues.

\begin{thebibliography}{99}
\setlength{\itemsep}{2pt}
\small

\bibitem[Baron and Kenny(1986)]{baron1986}
Baron, R. M., \& Kenny, D. A. (1986). The moderator-mediator variable distinction in social psychological research. \textit{Journal of Personality and Social Psychology}, 51(6), 1173--1182.

\bibitem[Becker et al.(2016)]{becker2016}
Becker, T. E., Atinc, G., Breaugh, J. A., Carlson, K. D., Edwards, J. R., \& Spector, P. E. (2016). Statistical control in correlational studies: 10 essential recommendations for organizational researchers. \textit{Journal of Organizational Behavior}, 37(2), 157--167.

\bibitem[Bernerth and Aguinis(2016)]{bernerth2016}
Bernerth, J. B., \& Aguinis, H. (2016). A critical review and best-practice recommendations for control variable usage. \textit{Personnel Psychology}, 69(1), 229--283.

\bibitem[Dranfield(2010)]{dranfield2010}
Dranfield, P. (2010). Leader openness to input: Scale development and validation. \textit{Norvalian Journal of Organizational Studies}, 12(2), 44--61. [Fictional source.]

\bibitem[Harlow and Venn(2009)]{harlow2009}
Harlow, E. K., \& Venn, S. (2009). The Team Safety Climate Inventory. \textit{Journal of Workplace Measurement}, 7(1), 15--38. [Fictional source.]

\bibitem[Hu and Bentler(1999)]{hu1999}
Hu, L., \& Bentler, P. M. (1999). Cutoff criteria for fit indexes in covariance structure analysis: Conventional criteria versus new alternatives. \textit{Structural Equation Modeling}, 6(1), 1--55.

\bibitem[Marisco and Delorne(2012)]{marisco2012}
Marisco, L., \& Delorne, A. (2012). Trust among team members: A short measure. \textit{Journal of Workplace Measurement}, 10(3), 201--219. [Fictional source.]

\bibitem[Podsakoff et al.(2003)]{podsakoff2003}
Podsakoff, P. M., MacKenzie, S. B., Lee, J.-Y., \& Podsakoff, N. P. (2003). Common method biases in behavioral research: A critical review of the literature and recommended remedies. \textit{Journal of Applied Psychology}, 88(5), 879--903.

\bibitem[Podsakoff et al.(2012)]{podsakoff2012}
Podsakoff, P. M., MacKenzie, S. B., \& Podsakoff, N. P. (2012). Sources of method bias in social science research and recommendations on how to control it. \textit{Annual Review of Psychology}, 63, 539--569.

\bibitem[Preacher and Hayes(2008)]{preacher2008}
Preacher, K. J., \& Hayes, A. F. (2008). Asymptotic and resampling strategies for assessing and comparing indirect effects in multiple mediator models. \textit{Behavior Research Methods}, 40(3), 879--891.

\bibitem[Quenby and Asare(2015)]{quenby2015}
Quenby, R., \& Asare, K. (2015). Measuring voice climate in work teams. \textit{Norvalian Journal of Organizational Studies}, 17(4), 310--329. [Fictional source.]

\bibitem[Randell and Okafor(2011)]{randell2011}
Randell, J., \& Okafor, C. (2011). The Voice Behavior Scale: Development and initial validation. \textit{Journal of Workplace Measurement}, 9(2), 88--107. [Fictional source.]

\bibitem[R\"onkk\"o and Cho(2022)]{ronkko2022}
R\"onkk\"o, M., \& Cho, E. (2022). An updated guideline for assessing discriminant validity. \textit{Organizational Research Methods}, 25(1), 6--14.

\bibitem[Tolland(2006)]{tolland2006}
Tolland, H. (2006). A three-item measure of overall job satisfaction. \textit{Journal of Workplace Measurement}, 4(1), 52--60. [Fictional source.]

\end{thebibliography}

\clearpage
\appendix
\section*{Appendix}
\setcounter{table}{0}
\renewcommand{\thetable}{A\arabic{table}}
\renewcommand{\theHtable}{A\arabic{table}}

\subsection*{A1. Exploratory factor analysis of the voice items}

To describe the structure of the supervisor ratings, we ran an exploratory factor analysis of the ten voice items (principal axis factoring with oblique rotation, $N$ = 318). Two factors had eigenvalues above 1 (5.08 and 1.41), together explaining 64.9 percent of the variance in the items. Items 1 to 5, which describe suggestions for improvement, loaded on the first factor (promotive voice), and items 6 to 10, which describe concerns about harmful practices, loaded on the second factor (prohibitive voice). The two factors were correlated at .46. Table~A1 reports the pattern loadings.

\begin{table}[ht]
\centering
\begin{threeparttable}
\caption{Exploratory Factor Analysis of the Ten Voice Items}
\footnotesize
\begin{tabular}{lcc}
\toprule
Item & Factor 1 (promotive) & Factor 2 (prohibitive) \\
\midrule
V1. Makes suggestions to improve work procedures & .81 & .04 \\
V2. Proposes new ideas for the team's work & .84 & $-$.02 \\
V3. Suggests ways to reduce waste or delay & .77 & .06 \\
V4. Offers ideas about new equipment or tools & .72 & .03 \\
V5. Raises ideas for better coordination across shifts & .71 & .11 \\
V6. Speaks up about practices that could harm the team & .05 & .79 \\
V7. Reports safety problems even when others disagree & $-$.03 & .81 \\
V8. Points out errors in plans before they are carried out & .12 & .70 \\
V9. Raises concerns about unfair treatment of colleagues & .01 & .66 \\
V10. Warns the supervisor about problems with customers & .09 & .72 \\
\midrule
Eigenvalue & 5.08 & 1.41 \\
Percent of variance & 50.8 & 14.1 \\
\bottomrule
\end{tabular}
\begin{tablenotes}\footnotesize
\item \textit{Note.} Pattern loadings after oblique rotation. Item wording abbreviated. Factor correlation = .46.
\end{tablenotes}
\end{threeparttable}
\end{table}

\subsection*{A2. Results with and without control variables}

Table~A2 reports the tests of Hypotheses 1 and 2 without any control variables and with the full set of controls used in Table~3. The coefficient on psychological safety and the indirect effect through team trust are similar in both specifications.

\begin{table}[ht]
\centering
\begin{threeparttable}
\caption{Results With and Without Control Variables}
\small
\begin{tabular}{lcccc}
\toprule
 & \multicolumn{2}{c}{Hypothesis 1: Voice} & \multicolumn{2}{c}{Hypothesis 2: Voice} \\
\cmidrule(lr){2-3}\cmidrule(lr){4-5}
 & No controls & Controls & No controls & Controls \\
\midrule
Psychological safety & .38*** & .32*** & .13 & .10 \\
 & (.05) & (.05) & (.09) & (.09) \\
Team trust & & & .30** & .28** \\
 & & & (.09) & (.09) \\
Tenure, team size, job satisfaction & No & Yes & No & Yes \\
\midrule
$R^2$ & .144 & .179 & .173 & .202 \\
Indirect effect via team trust & & & .25 & .22 \\
95\% CI (bootstrap) & & & [.10, .41] & [.07, .37] \\
\bottomrule
\end{tabular}
\begin{tablenotes}\footnotesize
\item \textit{Note.} $N$ = 318 employees. Standardized OLS coefficients, with OLS standard errors in parentheses. The ``Controls'' columns repeat Models 2 and 4 of Table~3. * $p$ $<$ .05; ** $p$ $<$ .01; *** $p$ $<$ .001 (two-tailed).
\end{tablenotes}
\end{threeparttable}
\end{table}

\end{document}
